\uuid{Upn1}
\exo7id{2425}
\titre{exo7 2425}
\auteur{matexo1}
\organisation{exo7}
\datecreate{2002-02-01}
\isIndication{false}
\isCorrection{false}
\chapitre{Espace vectoriel}
\sousChapitre{Définition, sous-espace}
\cpge{mpsi-mp2i.espaces-vectoriels.b, pcsi.espaces-vectoriels.b, ptsi.espaces-vectoriels.a, mpsi-mp2i.espaces-vectoriels.d, pcsi.espaces-vectoriels.d, ptsi.espaces-vectoriels.d, mpsi-mp2i.espaces-vectoriels.a, pcsi.espaces-vectoriels.a, ptsi.geometrie-espace.c, bcpst1.espace-vectoriel-kn.a}
\module{Algèbre}
\niveau{L1}
\difficulte{}

\contenu{
\texte{
\begin{center}
{\sc VRAI ou FAUX}
\end{center}
}
\begin{enumerate}
    \item \question{L'ensemble $\{0\}$ est un espace vectoriel r\'eel.}
    \item \question{L'ensemble $\{0,1\}$ est un espace vectoriel r\'eel.}
    \item \question{Tout sous-espace vectoriel autre que $\{0\}$ poss\`ede un
sous-espace strict.}
    \item \question{L'intersection de deux sous-espaces vectoriels (d'un m\^eme
espace plus grand) est un espace vectoriel.}
    \item \question{La r\'eunion de deux sous-espaces vectoriels est un
espace vectoriel.}
    \item \question{La somme de deux sous-espaces vectoriels est un espace
vectoriel.}
    \item \question{Le produit cart\'esien $E\times F$ de deux espaces vectoriels
est un espace vectoriel.}
\end{enumerate}
}
